\chapter{ساختمان داده‌ها}

\index{data structure@ساختمان داده}

یک \key{ساختمان داده} روشی برای ذخیره داده‌ها در حافظه کامپیوتر است. انتخاب ساختمان داده مناسب برای هر مسئله اهمیت بالایی دارد، زیرا هر ساختمان داده مزایا و معایب خاص خود را دارد. پرسش اساسی این است: کدام عملیات‌ها در ساختمان داده انتخاب‌شده کارآمد هستند؟

این فصل مهم‌ترین ساختمان داده‌های موجود در کتابخانه استاندارد C++ را معرفی می‌کند. استفاده از کتابخانه استاندارد تا جای ممکن توصیه می‌شود، زیرا زمان زیادی را صرفه‌جویی می‌کند. در ادامه کتاب، ساختمان داده‌های پیشرفته‌تری را بررسی خواهیم کرد که در کتابخانه استاندارد در دسترس نیستند.

\section{آرایه‌های پویا}

\index{dynamic array@آرایه پویا}
\index{vector@بردار (vector)}

یک \key{آرایه پویا} آرایه‌ای است که اندازه آن در طول اجرای برنامه قابل تغییر است. محبوب‌ترین آرایه پویا در C++ ساختار \texttt{vector} است که می‌توان از آن تقریباً مانند یک آرایه معمولی استفاده کرد.

کد زیر یک وکتور خالی ایجاد کرده و سه عنصر به آن اضافه می‌کند:

\begin{lstlisting}
vector<int> v;
v.push_back(3); // [3]
v.push_back(2); // [3,2]
v.push_back(5); // [3,2,5]
\end{lstlisting}

پس از این، دسترسی به عناصر مانند یک آرایه معمولی امکان‌پذیر است:

\begin{lstlisting}
cout << v[0] << "\n"; // 3
cout << v[1] << "\n"; // 2
cout << v[2] << "\n"; // 5
\end{lstlisting}

تابع \texttt{size} تعداد عناصر موجود در وکتور را برمی‌گرداند. کد زیر روی وکتور پیمایش کرده و تمام عناصر آن را چاپ می‌کند:

\begin{lstlisting}
for (int i = 0; i < v.size(); i++) {
    cout << v[i] << "\n";
}
\end{lstlisting}

\begin{samepage}
روش کوتاه‌تر برای پیمایش وکتور به صورت زیر است:

\begin{lstlisting}
for (auto x : v) {
    cout << x << "\n";
}
\end{lstlisting}
\end{samepage}

تابع \texttt{back} آخرین عنصر وکتور را برمی‌گرداند و تابع \texttt{pop\_back} آخرین عنصر را حذف می‌کند:

\begin{lstlisting}
vector<int> v;
v.push_back(5);
v.push_back(2);
cout << v.back() << "\n"; // 2
v.pop_back();
cout << v.back() << "\n"; // 5
\end{lstlisting}

کد زیر یک وکتور با پنج عنصر ایجاد می‌کند:

\begin{lstlisting}
vector<int> v = {2,4,2,5,1};
\end{lstlisting}

روش دیگر ساخت وکتور، تعیین تعداد عناصر و مقدار اولیه هر عنصر است:

\begin{lstlisting}
// size 10, initial value 0
vector<int> v(10);
\end{lstlisting}
\begin{lstlisting}
// size 10, initial value 5
vector<int> v(10, 5);
\end{lstlisting}

پیاده‌سازی داخلی وکتور از یک آرایه معمولی استفاده می‌کند. اگر اندازه وکتور افزایش یابد و آرایه بیش از حد کوچک شود، آرایه جدیدی تخصیص داده شده و تمام عناصر به آرایه جدید منتقل می‌شوند. با این حال، این اتفاق مکرراً رخ نمی‌دهد و پیچیدگی زمانی میانگین \texttt{push\_back} برابر با $O(1)$ است.

\index{string@رشته (string)}

ساختار \texttt{string}
هم آرایه پویا است، تقریباً مثل بردار کار می‌کند.
رشته‌ها نحو اختصاصی دارند
که در سایر ساختارهای داده نیست.
الحاق رشته‌ها با نماد \texttt{+} انجام می‌شود.
تابع $\texttt{substr}(k,x)$ زیررشته
از موقعیت $k$ با طول $x$ برمی‌گرداند،
تابع $\texttt{find}(\texttt{t})$ موقعیت
اولین وقوع زیررشته \texttt{t} پیدا می‌کند.

کد زیر عملیات رشته را نشان می‌دهد:

\begin{lstlisting}
string a = "hatti";
string b = a+a;
cout << b << "\n"; // hattihatti
b[5] = 'v';
cout << b << "\n"; // hattivatti
string c = b.substr(3,4);
cout << c << "\n"; // tiva
\end{lstlisting}

\section{ساختارهای مجموعه}

\index{مجموعه}

یک \key{مجموعه} داده‌ساختاری است که
دسته‌ای از عناصر را نگه می‌دارد.
عملیات اصلی:
درج، جستجو، حذف عنصر.

کتابخانه استاندارد C++ دو پیاده‌سازی دارد:
ساختار \texttt{set} بر پایه درخت دودویی
متوازن است، عملیات آن در زمان $O(\log n)$ اجرا می‌شود.
ساختار \texttt{unordered\_set} از درهم‌سازی استفاده می‌کند،
عملیات آن در حالت میانگین در زمان $O(1)$ اجرا می‌شود.

انتخاب نوع پیاده‌سازی سلیقه‌ای است.
مزیت ساختار \texttt{set}:
حفظ ترتیب عناصر،
ارائه توابعی که
در \texttt{unordered\_set} نیست.
در مقابل، \texttt{unordered\_set}
می‌تواند کارآمدتر باشد.

کد زیر مجموعه اعداد صحیح
می‌سازد،
برخی عملیات را نشان می‌دهد.
تابع \texttt{insert} عنصر به مجموعه اضافه می‌کند،
تابع \texttt{count} تعداد تکرار
عنصر در مجموعه را برمی‌گرداند،
تابع \texttt{erase} عنصر از مجموعه حذف می‌کند.

\begin{lstlisting}
set<int> s;
s.insert(3);
s.insert(2);
s.insert(5);
cout << s.count(3) << "\n"; // 1
cout << s.count(4) << "\n"; // 0
s.erase(3);
s.insert(4);
cout << s.count(3) << "\n"; // 0
cout << s.count(4) << "\n"; // 1
\end{lstlisting}

مجموعه مثل بردار استفاده می‌شود،
اما دسترسی به عناصر
با نماد \texttt{[]} ممکن نیست.
کد زیر مجموعه می‌سازد،
تعداد عناصر را چاپ می‌کند، سپس
روی همه عناصر پیمایش می‌کند:
\begin{lstlisting}
set<int> s = {2,5,6,8};
cout << s.size() << "\n"; // 4
for (auto x : s) {
    cout << x << "\n";
}
\end{lstlisting}

ویژگی مهم مجموعه:
تمام عناصر \emph{متمایز} هستند.
بنابراین تابع \texttt{count} همیشه
یا 0 (عنصر نیست)
یا 1 (عنصر هست) برمی‌گرداند،
تابع \texttt{insert} عنصری که قبلاً باشد را
دوباره اضافه نمی‌کند.
کد زیر این رفتار را نشان می‌دهد:

\begin{lstlisting}
set<int> s;
s.insert(5);
s.insert(5);
s.insert(5);
cout << s.count(5) << "\n"; // 1
\end{lstlisting}

سی‌پلاس‌پلاس همچنین شامل ساختارهای
\texttt{multiset} و \texttt{unordered\_multiset}
است که عملکردی مشابه \texttt{set}
و \texttt{unordered\_set} دارند،
اما می‌توانند شامل چندین نسخه از یک عنصر باشند.
به عنوان مثال، در کد زیر هر سه نسخه از
عدد ۵ به یک چندمجموعه اضافه می‌شوند:

\begin{lstlisting}
multiset<int> s;
s.insert(5);
s.insert(5);
s.insert(5);
cout << s.count(5) << "\n"; // 3
\end{lstlisting}
تابع \texttt{erase} تمام نسخه‌های یک عنصر را
از چندمجموعه حذف می‌کند:
\begin{lstlisting}
s.erase(5);
cout << s.count(5) << "\n"; // 0
\end{lstlisting}
اغلب، تنها حذف یک نسخه از عنصر مورد نظر است
که می‌تواند به صورت زیر انجام شود:
\begin{lstlisting}
s.erase(s.find(5));
cout << s.count(5) << "\n"; // 2
\end{lstlisting}

\section{ساختارهای نگاشت (Map)}

\index{نگاشت}

یک \key{نگاشت} یک آرایه تعمیم‌یافته است
که از جفت‌های کلید-مقدار تشکیل می‌شود.
در حالی که کلیدها در یک آرایه معمولی همیشه
اعداد صحیح متوالی $0,1,\ldots,n-1$ هستند
(که $n$ اندازه آرایه است)،
کلیدها در یک نگاشت می‌توانند از هر نوع داده‌ای باشند و
نیازی نیست که مقادیری متوالی باشند.

کتابخانه استاندارد C++ شامل دو پیاده‌سازی نگاشت است
که متناظر با پیاده‌سازی‌های مجموعه‌اند: ساختار
\texttt{map} بر پایه یک درخت دودویی متوازن
است و دسترسی به عناصر
زمان $O(\log n)$ می‌برد،
در حالی که ساختار
\texttt{unordered\_map} از درهم‌سازی استفاده می‌کند
و دسترسی به عناصر در حالت میانگین زمان $O(1)$ می‌برد.

کد زیر یک نگاشت ایجاد می‌کند
که در آن کلیدها رشته و مقادیر اعداد صحیح هستند:

\begin{lstlisting}
map<string,int> m;
m["monkey"] = 4;
m["banana"] = 3;
m["harpsichord"] = 9;
cout << m["banana"] << "\n"; // 3
\end{lstlisting}

اگر مقدار یک کلید درخواست شود
اما نگاشت شامل آن نباشد،
کلید به طور خودکار با یک مقدار پیش‌فرض
به نگاشت اضافه می‌شود.
برای مثال در کد زیر،
کلید ''aybabtu'' با مقدار 0
به نگاشت اضافه می‌شود.

\begin{lstlisting}
map<string,int> m;
cout << m["aybabtu"] << "\n"; // 0
\end{lstlisting}
تابع \texttt{count} بررسی می‌کند
که آیا یک کلید در نگاشت وجود دارد یا خیر:
\begin{lstlisting}
if (m.count("aybabtu")) {
    // key exists
}
\end{lstlisting}
کد زیر تمام کلیدها و مقادیر
موجود در یک نگاشت را چاپ می‌کند:
\begin{lstlisting}
for (auto x : m) {
    cout << x.first << " " << x.second << "\n";
}
\end{lstlisting}

\section{تکرارکننده‌ها و بازه‌ها}

\index{تکرارکننده}

بسیاری از توابع در کتابخانه استاندارد C++
با تکرارکننده‌ها کار می‌کنند.
یک \key{تکرارکننده} متغیری است که به یک عنصر
در یک ساختار داده اشاره می‌کند.

تکرارکننده‌های پرکاربرد \texttt{begin}
و \texttt{end} بازه‌ای را تعریف می‌کنند که شامل
تمام عناصر موجود در یک ساختار داده است.
تکرارکننده \texttt{begin} به
اولین عنصر در ساختار داده اشاره می‌کند،
و تکرارکننده \texttt{end} به
موقعیت \emph{بعد} از آخرین عنصر اشاره دارد.
وضعیت به شکل زیر است:

\begin{center}
\begin{tabular}{llllllllll}
\{ & 3, & 4, & 6, & 8, & 12, & 13, & 14, & 17 & \} \\
& $\uparrow$ & & & & & & & & $\uparrow$ \\
& \multicolumn{3}{l}{\texttt{s.begin()}} & & & & & & \texttt{s.end()} \\
\end{tabular}
\end{center}

به عدم تقارن در تکرارکننده‌ها دقت کنید:
\texttt{s.begin()} به عنصری درون ساختمان داده اشاره دارد،
در حالی که \texttt{s.end()} به خارج از ساختمان داده اشاره می‌کند.
بنابراین، بازه تعریف‌شده توسط تکرارکننده‌ها \emph{نیمه‌باز} است.

\subsubsection{کار با بازه‌ها}

تکرارکننده‌ها در توابع کتابخانه استاندارد ++C کاربرد دارند،
توابعی که بازه‌ای از عناصر یک ساختمان داده به آن‌ها داده می‌شود.
معمولاً پردازش همه عناصر یک ساختمان داده مد نظر است،
بنابراین تکرارکننده‌های \texttt{begin} و \texttt{end} به تابع ورودی داده می‌شوند.

برای نمونه، کد زیر یک بردار را با استفاده از تابع \texttt{sort} مرتب می‌کند،
سپس ترتیب عناصر را با تابع \texttt{reverse} معکوس می‌سازد،
و در نهایت ترتیب عناصر را به کمک تابع \texttt{random\_shuffle} برمی‌زند.

\index{sort@\texttt{sort}}
\index{reverse@\texttt{reverse}}
\index{random\_shuffle@\texttt{random\_shuffle}}

\begin{lstlisting}
sort(v.begin(), v.end());
reverse(v.begin(), v.end());
random_shuffle(v.begin(), v.end());
\end{lstlisting}

این توابع برای آرایه‌های معمولی نیز کاربرد دارند.
در این حالت، به جای تکرارکننده، اشاره‌گرهایی به آرایه ورودی داده می‌شوند:

\newpage
\begin{lstlisting}
sort(a, a+n);
reverse(a, a+n);
random_shuffle(a, a+n);
\end{lstlisting}

\subsubsection{تکرارکننده‌های مجموعه}

تکرارکننده‌ها اغلب برای دسترسی به عناصر یک مجموعه استفاده می‌شوند.
کد زیر تکرارکننده \texttt{it} را ایجاد می‌کند که به کوچک‌ترین عنصر مجموعه اشاره دارد:
\begin{lstlisting}
set<int>::iterator it = s.begin();
\end{lstlisting}
روش کوتاه‌تر برای نوشتن این کد:
\begin{lstlisting}
auto it = s.begin();
\end{lstlisting}
عنصری که تکرارکننده به آن اشاره می‌کند با نماد \texttt{*} قابل دسترسی است.
برای نمونه، کد زیر نخستین عنصر مجموعه را چاپ می‌کند:

\begin{lstlisting}
auto it = s.begin();
cout << *it << "\n";
\end{lstlisting}

تکرارکننده‌ها را می‌توان با عملگرهای \texttt{++} (به جلو) و \texttt{--} (به عقب) حرکت داد،
بدین معنی که تکرارکننده به عنصر بعدی یا قبلی درون مجموعه جابه‌جا می‌شود.

کد زیر تمام عناصر را به ترتیب صعودی چاپ می‌کند:
\begin{lstlisting}
for (auto it = s.begin(); it != s.end(); it++) {
    cout << *it << "\n";
}
\end{lstlisting}
کد زیر بزرگ‌ترین عنصر مجموعه را چاپ می‌کند:
\begin{lstlisting}
auto it = s.end(); it--;
cout << *it << "\n";
\end{lstlisting}

تابع $\texttt{find}(x)$ تکرارکننده‌ای برمی‌گرداند که به عنصری با مقدار $x$ اشاره دارد.
با این حال، اگر مجموعه شامل $x$ نباشد، خروجی تکرارکننده \texttt{end} خواهد بود.

\begin{lstlisting}
auto it = s.find(x);
if (it == s.end()) {
    // x پیدا نشد
}
\end{lstlisting}

تابع $\texttt{lower\_bound}(x)$ یک تکرارکننده به کوچک‌ترین عنصر مجموعه برمی‌گرداند که مقدار آن \emph{حداقل} $x$ باشد، و تابع $\texttt{upper\_bound}(x)$ یک تکرارکننده به کوچک‌ترین عنصر مجموعه برمی‌گرداند که مقدار آن \emph{بزرگ‌تر از} $x$ باشد. در هر دو تابع، اگر چنین عنصری وجود نداشته باشد، مقدار بازگشتی \texttt{end} است. این توابع در ساختار داده \texttt{unordered\_set} که ترتیب عناصر را حفظ نمی‌کند، پشتیبانی نمی‌شوند.

\begin{samepage}
برای نمونه، کد زیر نزدیک‌ترین عنصر به $x$ را می‌یابد:

\begin{lstlisting}
auto it = s.lower_bound(x);
if (it == s.begin()) {
    cout << *it << "\n";
} else if (it == s.end()) {
    it--;
    cout << *it << "\n";
} else {
    int a = *it; it--;
    int b = *it;
    if (x-b < a-x) cout << b << "\n";
    else cout << a << "\n";
}
\end{lstlisting}

این کد فرض می‌کند که مجموعه خالی نیست، و با استفاده از تکرارکننده \texttt{it} تمام حالت‌های ممکن را بررسی می‌کند. ابتدا، تکرارکننده به کوچک‌ترین عنصری اشاره می‌کند که مقدار آن حداقل $x$ است. اگر \texttt{it} برابر با \texttt{begin} باشد، عنصر متناظر نزدیک‌ترین مقدار به $x$ است. اگر \texttt{it} برابر با \texttt{end} باشد، بزرگ‌ترین عنصر مجموعه نزدیک‌ترین مقدار به $x$ است. اگر هیچ‌یک از حالت‌های قبلی رخ ندهد، عنصر نزدیک‌تر به $x$ یا عنصر متناظر با \texttt{it} است یا عنصر پیش از آن.
\end{samepage}

\section{سایر ساختارها}

\subsubsection{بیت‌ست (Bitset)}

\index{bitset}

یک \key{بیت‌ست} آرایه‌ای است که هر مقدار آن ۰ یا ۱ است. برای نمونه، کد زیر یک بیت‌ست با ۱۰ عنصر می‌سازد:
\begin{lstlisting}
bitset<10> s;
s[1] = 1;
s[3] = 1;
s[4] = 1;
s[7] = 1;
cout << s[4] << "\n"; // 1
cout << s[5] << "\n"; // 0
\end{lstlisting}

مزیت استفاده از بیت‌ست‌ها مصرف حافظه کمتر آن‌ها نسبت به آرایه‌های معمولی است، زیرا هر عنصر در بیت‌ست تنها یک بیت حافظه می‌گیرد. برای نمونه، اگر $n$ بیت درون یک آرایه \texttt{int} ذخیره شود، $32n$ بیت حافظه مصرف خواهد شد، اما بیت‌ست متناظر تنها نیازمند $n$ بیت حافظه است. افزون بر این، مقادیر بیت‌ست را می‌توان به‌شکل کارآمد با عملگرهای بیتی دستکاری کرد، که این کار بهینه‌سازی الگوریتم‌ها به کمک بیت‌ست را ممکن می‌سازد.

کد زیر روش دیگری برای ساخت بیت‌ست بالا نشان می‌دهد:
\begin{lstlisting}
bitset<10> s(string("0010011010")); // from right to left
cout << s[4] << "\n"; // 1
cout << s[5] << "\n"; // 0
\end{lstlisting}

تابع \texttt{count} تعداد یک‌ها را در بیت‌ست برمی‌گرداند:

\begin{lstlisting}
bitset<10> s(string("0010011010"));
cout << s.count() << "\n"; // 4
\end{lstlisting}

کد زیر نمونه‌هایی از به‌کارگیری عملگرهای بیتی را نشان می‌دهد:
\begin{lstlisting}
bitset<10> a(string("0010110110"));
bitset<10> b(string("1011011000"));
cout << (a&b) << "\n"; // 0010010000
cout << (a|b) << "\n"; // 1011111110
cout << (a^b) << "\n"; // 1001101110
\end{lstlisting}

\subsubsection{دک (Deque)}

\index{deque}

یک \key{deque} (دک یا صف دوطرفه) آرایه‌ای پویا است
که اندازه آن می‌تواند به شکلی کارآمد
از هر دو سر آرایه تغییر کند.
همانند بردار، یک deque توابع
\texttt{push\_back} و \texttt{pop\_back} را ارائه می‌دهد، اما
همچنین شامل توابع
\texttt{push\_front} و \texttt{pop\_front} است
که در بردار وجود ندارند.

یک deque می‌تواند به صورت زیر استفاده شود:
\begin{lstlisting}
deque<int> d;
d.push_back(5); // [5]
d.push_back(2); // [5,2]
d.push_front(3); // [3,5,2]
d.pop_back(); // [3,5]
d.pop_front(); // [5]
\end{lstlisting}

پیاده‌سازی داخلی یک deque
پیچیده‌تر از یک بردار است،
و به همین دلیل، deque کندتر از بردار است.
با این حال، هر دو عملیات افزودن و حذف
عناصر در هر دو سر به طور میانگین به زمان $O(1)$ نیاز دارند.

\subsubsection{پشته}

\index{stack}

یک \key{پشته} (stack)
ساختمان داده‌ای است که دو عملیات
با زمان $O(1)$ ارائه می‌دهد:
افزودن یک عنصر به بالا،
و حذف یک عنصر از بالا.
تنها دسترسی به عنصر
بالای پشته امکان‌پذیر است.

کد زیر نحوه استفاده از پشته را نشان می‌دهد:
\begin{lstlisting}
stack<int> s;
s.push(3);
s.push(2);
s.push(5);
cout << s.top(); // 5
s.pop();
cout << s.top(); // 2
\end{lstlisting}
\subsubsection{صف}

\index{queue}

یک \key{صف} (queue) نیز
دو عملیات با زمان $O(1)$ ارائه می‌دهد:
افزودن یک عنصر به انتهای صف،
و حذف عنصر اول از صف.
تنها دسترسی به عناصر اول
و آخر صف امکان‌پذیر است.

کد زیر نحوه استفاده از صف را نشان می‌دهد:
\begin{lstlisting}
queue<int> q;
q.push(3);
q.push(2);
q.push(5);
cout << q.front(); // 3
q.pop();
cout << q.front(); // 2
\end{lstlisting}

\subsubsection{صف اولویت‌دار}

\index{priority queue}
\index{heap}

یک \key{صف اولویت‌دار} (priority queue)
مجموعه‌ای از عناصر را نگهداری می‌کند.
عملیات پشتیبانی‌شده شامل درج و،
بسته به نوع صف،
بازیابی و حذف
کوچک‌ترین یا بزرگ‌ترین عنصر است.
درج و حذف به زمان $O(\log n)$ نیاز دارند،
و بازیابی در زمان $O(1)$ انجام می‌شود.

در حالی که یک مجموعه مرتب تمام عملیات
یک صف اولویت‌دار را به صورت کارآمد پشتیبانی می‌کند،
مزیت استفاده از صف اولویت‌دار در این است
که ضرایب ثابت کوچک‌تری دارد.
یک صف اولویت‌دار معمولاً با استفاده از
ساختار هرم پیاده‌سازی می‌شود که بسیار ساده‌تر از
درخت دودویی متعادل مورد استفاده در مجموعه مرتب است.

\begin{samepage}
به طور پیش‌فرض، عناصر در یک صف اولویت‌دار
در ++C به صورت نزولی مرتب شده‌اند،
و یافتن و حذف
بزرگ‌ترین عنصر در صف ممکن است.
کد زیر این موضوع را نشان می‌دهد:

\begin{lstlisting}
priority_queue<int> q;
q.push(3);
q.push(5);
q.push(7);
q.push(2);
cout << q.top() << "\n"; // 7
q.pop();
cout << q.top() << "\n"; // 5
q.pop();
q.push(6);
cout << q.top() << "\n"; // 6
q.pop();
\end{lstlisting}
\end{samepage}

اگر بخواهیم یک صف اولویت‌دار ایجاد کنیم
که از یافتن و حذف
کوچک‌ترین عنصر پشتیبانی کند،
می‌توانیم به صورت زیر عمل کنیم:

\begin{lstlisting}
priority_queue<int,vector<int>,greater<int>> q;
\end{lstlisting}

\subsubsection{ساختمان داده‌های مبتنی بر سیاست}

کامپایلر \texttt{g++} از برخی ساختارهای داده که بخشی از کتابخانه استاندارد ++C نیستند نیز پشتیبانی می‌کند.
این ساختارها \emph{مبتنی بر خط‌مشی} (policy-based) نامیده می‌شوند.
برای استفاده از این ساختارها، خطوط زیر باید به کد اضافه شوند:
\begin{lstlisting}
#include <ext/pb_ds/assoc_container.hpp>
using namespace __gnu_pbds; 
\end{lstlisting}
سپس می‌توان ساختار داده \texttt{indexed\_set} را تعریف کرد که مانند \texttt{set} است ولی مثل آرایه اندیس‌پذیر است.
تعریف برای مقادیر \texttt{int}:
\begin{lstlisting}
typedef tree<int,null_type,less<int>,rb_tree_tag,
             tree_order_statistics_node_update> indexed_set; 
\end{lstlisting}
ساخت مجموعه:
\begin{lstlisting}
indexed_set s;
s.insert(2);
s.insert(3);
s.insert(7);
s.insert(9);
\end{lstlisting}
ویژگی این مجموعه: دسترسی به اندیس‌هایی که عناصر در آرایه مرتب‌شده می‌داشتند.
تابع $\texttt{find\_by\_order}$ تکرارکننده (iterator) به عنصر در موقعیت داده‌شده برمی‌گرداند:
\begin{lstlisting}
auto x = s.find_by_order(2);
cout << *x << "\n"; // 7
\end{lstlisting}
تابع $\texttt{order\_of\_key}$ موقعیت عنصر داده‌شده را برمی‌گرداند:
\begin{lstlisting}
cout << s.order_of_key(7) << "\n"; // 2
\end{lstlisting}
اگر عنصر در مجموعه نباشد، موقعیتی که عنصر در صورت حضور می‌داشت برمی‌گردد:
\begin{lstlisting}
cout << s.order_of_key(6) << "\n"; // 2
cout << s.order_of_key(8) << "\n"; // 3
\end{lstlisting}
هر دو تابع در زمان لگاریتمی اجرا می‌شوند.

\section{مقایسه با مرتب‌سازی}

حل مسئله با استفاده از ساختار داده‌ها یا مرتب‌سازی معمولاً ممکن است.
گاهی تفاوت‌های چشمگیری در کارایی واقعی این روش‌ها وجود دارد که در پیچیدگی زمانی آن‌ها پنهان است.

مسئله: دو لیست $A$ و $B$ با $n$ عنصر داده شده است.
هدف: شمارش تعداد عناصری که در هر دو لیست حضور دارند.
مثال برای لیست‌های
\[A = [5,2,8,9] \hspace{10px} \textrm{و} \hspace{10px} B = [3,2,9,5],\]
پاسخ ۳ است زیرا اعداد ۲، ۵ و ۹ در هر دو لیست هستند.

راه‌حل مستقیم: بررسی تمام جفت‌های عناصر در زمان $O(n^2)$.
الگوریتم‌های کارآمدتر:

\subsubsection{الگوریتم ۱}

ساخت مجموعه از عناصر $A$، سپس پیمایش عناصر $B$ و بررسی حضور هر عنصر در $A$.
کارآمد است چون عناصر $A$ داخل یک مجموعه‌اند.
با ساختار \texttt{set}، پیچیدگی زمانی الگوریتم $O(n \log n)$ است.

\subsubsection{الگوریتم ۲}

حفظ ترتیب عناصر لازم نیست، پس به جای \texttt{set} می‌توان از \texttt{unordered\_set} استفاده کرد.
روش ساده برای افزایش کارایی چون فقط ساختار داده تغییر می‌کند.
پیچیدگی زمانی الگوریتم جدید $O(n)$ است.

\subsubsection{الگوریتم ۳}

به جای ساختمان داده‌ها، می‌توان از مرتب‌سازی استفاده کرد.
ابتدا، هر دو لیست $A$ و $B$ را مرتب می‌کنیم.
سپس، هم‌زمان روی هر دو لیست پیمایش کرده
و عناصر مشترک را پیدا می‌کنیم.
پیچیدگی زمانی مرتب‌سازی $O(n \log n)$ است
و باقی الگوریتم در زمان $O(n)$ اجرا می‌شود،
بنابراین پیچیدگی زمانی کل $O(n \log n)$ است.

\subsubsection{مقایسه کارایی}

جدول زیر میزان کارایی الگوریتم‌های فوق را
به ازای مقادیر مختلف $n$ نشان می‌دهد،
در حالی که عناصر لیست‌ها اعداد صحیح تصادفی
بین $1 \ldots 10^9$ هستند:

\begin{center}
\begin{tabular}{rrrr}
$n$ & الگوریتم ۱ & الگوریتم ۲ & الگوریتم ۳ \\
\hline
$10^6$ & $1.5$ ثانیه & $0.3$ ثانیه & $0.2$ ثانیه \\
$2 \cdot 10^6$ & $3.7$ ثانیه & $0.8$ ثانیه & $0.3$ ثانیه \\
$3 \cdot 10^6$ & $5.7$ ثانیه & $1.3$ ثانیه & $0.5$ ثانیه \\
$4 \cdot 10^6$ & $7.7$ ثانیه & $1.7$ ثانیه & $0.7$ ثانیه \\
$5 \cdot 10^6$ & $10.0$ ثانیه & $2.3$ ثانیه & $0.9$ ثانیه \\
\end{tabular}
\end{center}

الگوریتم‌های ۱ و ۲ یکسان هستند به جز اینکه
از ساختارهای مجموعه متفاوتی استفاده می‌کنند.
در این مسئله، این انتخاب تاثیر مهمی روی
زمان اجرا دارد، زیرا الگوریتم ۲
بین ۴ تا ۵ برابر سریع‌تر از الگوریتم ۱ است.

با این حال، کارآمدترین الگوریتم، الگوریتم ۳ است
که از مرتب‌سازی استفاده می‌کند.
این الگوریتم تنها نیمی از زمان الگوریتم ۲ را مصرف می‌کند.
نکته جالب این است که پیچیدگی زمانی هر دو
الگوریتم ۱ و الگوریتم ۳ برابر با $O(n \log n)$ است،
اما با این وجود، الگوریتم ۳ ده برابر سریع‌تر است.
دلیل این موضوع ساده بودن فرآیند مرتب‌سازی است
که تنها یک بار در ابتدای الگوریتم ۳ انجام می‌شود
و باقی الگوریتم در زمان خطی کار می‌کند.
در مقابل،
الگوریتم ۱ یک درخت دودویی متوازن پیچیده را
در تمام طول اجرای الگوریتم نگهداری می‌کند.